\ifx\danslelivre\undefined
\RequirePackage{luatex85}
\documentclass[10pt]{../fichiers-configuration-latex/smfart}
\usepackage{polyglossia}
%\usepackage[francais]{babel}
\usepackage{fontspec}
\usepackage[math-style=ISO,bold-style=upright,sans-style=upright]{unicode-math}
%\usepackage[mathletters]{ucs}
%\usepackage[utf8]{inputenc}
%\usepackage[T1]{fontenc}
\usepackage[unicode,colorlinks=true,linkcolor=blue,citecolor=Sepia,urlcolor=magenta]{hyperref}
\setmainlanguage{french}
\setmainfont[Ligatures=TeX]{Linux Libertine O}
% variante possible : TeXGyrePagella ?
\setmathfont{XITS Math}

% maths
\usepackage{amsmath}
%\usepackage{amssymb}
%\usepackage{mathrsfs}
%\usepackage{stmaryrd}
%\usepackage{euscript}

% couleurs
\usepackage[usenames,dvipsnames]{xcolor}

% graphiques
\usepackage{graphicx}
\usepackage[all]{xy}
% Xy-pic 3.8.9's PDF driver assumes that \fontdimen8\textfont3 is a
% legacy TeX rule width.  Under LuaTeX with unicode-math/XITS Math it is
% instead 6.62pt, which turns every arrow shaft into a thick black bar.
% Override only xypdf's private cached width; do not alter the math font.
\makeatletter
\AtBeginDocument{%
  \@ifundefined{xP@preclw}{}{%
    \edef\xP@preclw{0.4pt}%
    \edef\xP@lw{\xP@dim\xP@preclw}%
    \PackageInfo{sga4-xypdf-compat}{Xy-pic PDF line width set to \xP@preclw}%
  }%
}
\makeatother
%\usepackage{tikz}
%\usetikzlibrary{matrix,arrows,calc}
%\tikzset{column sep=3em,row sep=2em,text height=1.5ex, text depth=0.25ex,>=angle 90}
%\tikzset{description/.style={fill=white,inner sep=1pt}}

% hyperliens, 'etiquettes
\usepackage{xr-hyper}
\usepackage{../fichiers-configuration-latex/smfthm}
%\ifx\cacheretiquettes\undefined
%\usepackage[inline]{showlabels} \else \fi % pour rendre visibles les 'etiquettes
%\usepackage{../fichiers-configuration-latex/makeidx} % pour faire un index
%\usepackage[pagebackref,unicode]{hyperref}  % pour voir o`u sont cit'ees les r'ef'erences
\hypersetup{unicode,colorlinks=true,linkcolor=blue,citecolor=Sepia,urlcolor=magenta,
  hypertexnames=false}
%\hypersetup{breaklinks=true} %permet le retour `a la ligne dans les liens trop longs. Incompatible avec pdflatex
%\renewcommand*{\backref}[1]{{\small [cit\'e page(s)~#1.]}}

\usepackage{xspace}
\usepackage{pifont} % pour expose XVII

\renewcommand{\thefootnote}{\roman{footnote}} % pour mettre les notes en bas de page sous la forme (ii) et non (2), qui ressemble trop (?) a un exposant.

\newcounter{pageoriginal}
\newcommand{\pageoriginale}{\refstepcounter{pageoriginal}\marginpar{\footnotesize\textbf{\thepageoriginal}}\xspace}
\let\pageoriginaled\pageoriginale

\usepackage{ifthen}
\newboolean{original}
\def\sisi{\ifthenelse{\boolean{original}}}
\setboolean{original}{false}

\newcounter{cnote}
\def\nde#1{\ifthenelse{\boolean{original}}{}{\stepcounter{cnote}\def\thefootnote{\arabic{cnote}}\protect\footnote[\value{cnote}]{N.D.E.: #1}}}
\def\ndemark{\ifthenelse{\boolean{original}}{}{\stepcounter{cnote}\def\thefootnote{\arabic{cnote}}\protect\footnotemark[\value{cnote}]}}
\def\ndetext#1{\ifthenelse{\boolean{original}}{}{\def\thefootnote{\arabic{cnote}}\protect\footnotetext[\value{cnote}]{N.D.E.: #1}}}

\setmainlanguage{french}
\setotherlanguage{english}

% unicode-math deliberately disables the legacy \Arrowvert glyph.  The
% source uses it once as a rotated equality marker; \Vert has the intended
% appearance and keeps the historical diagram buildable with LuaLaTeX.
\renewcommand{\Arrowvert}{\Vert}


%%
%% Pour afficher le numéro de version Git (exige que LuaTeX soit appelé avec --shell-escape)
%%

\catcode`\%=12
\directlua{
function command_output(cmd)
  local output_pipe = io.popen(cmd)
  local line_iter = output_pipe and output_pipe:lines()
  local line = line_iter and line_iter()
  if output_pipe then output_pipe:close() end
  return line
end
local git_command1 = [[git log --pretty=format:"%h %ai" -1]]
local git_command2 = [[if git diff --quiet HEAD ; then echo clean ; else echo dirty ; fi]]
function git_status()
  local status = command_output(git_command1)
  if status then
    tex.tprint({luatexbase.catcodetables.CatcodeTableLaTeX, string.char(0x5c).."texttt{"}, {-2,status}, {luatexbase.catcodetables.CatcodeTableLaTeX, "}"})
    if not (command_output(git_command2) == "clean") then tex.sprint(-2, " (modifications non validées)") end
  else
    tex.sprint(-2, "?")
  end
end
}
\catcode`\%=14
\def\showgitstatus{\directlua{git_status()}}


% macros pour LaTeX
\renewcommand{\mathsf}{}

% Police différente pour U+2113 SCRIPT SMALL L
\setmathfont[range="2113]{latinmodern-math.otf}

% Give every explicit SGA label its own destination and preserve the printed
% number supplied as the optional argument.  The former definition inherited
% whichever counter anchor happened to precede it and ignored #1 entirely.
\makeatletter
\def\etiquette[#1]#2{%
  \phantomsection
  \protected@edef\@currentlabel{#1}%
  \label{#2}%
}
\makeatother
\def\eqhyperref[#1]#2{(\hyperref[#1]{#2})}
\def\pagehyperref[#1]#2{(page \hyperref[#1]{#2})}
\def\hyperrefext[#1][#2]#3{\hyperref[#2]{#1 #3}}
\def\eqhyperrefext[#1][#2]#3{\hyperref[#2]{#1 #3}}     %cf. exposes 17
\def\neweqhyperrefext[#1][#2]#3{\hyperref[#2]{#1 #3}}  %cf. exposes 17

\def\shat#1{\protect\widehat{#1}}
\def\stilde#1{\protect\widetilde{#1}}
%%%%%%%%%%%%%%%%%%


\def\etcn{etc.\xspace}
\def\etco{etc...\xspace}
\def\etc{\sisi{\etco}{\etcn}}
\def\ldotsn{\sisi{...}{\etcn}}

\def\cf{{cf.}\kern.3em}
\def\Cf{{Cf.}\kern.3em}
\def\ie{{i.e.}\kern.3em}
\def\iev{{i.e.},\xspace}
\def\loccito{loc.\kern2pt cit.\xspace}
\def\loccitn{\textit{loc.\kern2pt cit.}\xspace}
\def\loccit{\sisi{\loccito}{\loccitn}}
\def\EGA{\textup{EGA}\kern.2em}
\def\SGA{\textup{SGA}\kern.3em}

\def\quoi{}

\def\quav{\quad,\quad}
\def\quaet{\quad \text{et}\quad}
\def\fibreproduct#1#2#3{#1 ×_{#3} #2}
\let\fprod\fibreproduct
\def\ccheck#1#2{\mbox{${\check{\check{\rm #1}}}^{\lower2pt\hbox{$\scriptstyle #2$}}$}}
\def\sheaf#1{\protect\mathcal{#1}}
\def\vieuxsheaf#1{{#1}} % cf. problème avec sheaf{a\Phi_i} par exemple...
\DeclareMathOperator{\Hom}{Hom}
\DeclareMathOperator{\can}{can}
\DeclareMathOperator{\Ind}{Ind}
\DeclareMathOperator{\ind}{ind}
%\DeclareMathOperator{\Hhomtop}{\mathcal{H}\!\mathit{omtop}}
\DeclareMathOperator{\Hhomtop}{\mathbf{Homtop}}
\DeclareMathOperator{\Point}{Point}
\DeclareMathOperator{\Pro}{Pro}
\DeclareMathOperator{\Top}{Top}
\DeclareMathOperator{\Mod}{Mod}
\DeclareMathOperator{\Cart}{Cart}
%\DeclareMathOperator{\Mmorsite}{\mathcal{M}\!\mathit{orsite}}
\DeclareMathOperator{\Mmorsite}{\mathbf{Morsite}}
\DeclareMathOperator{\loc}{loc}
\DeclareMathOperator{\Loc}{Loc}
\DeclareMathOperator{\Pic}{Pic}
\DeclareMathOperator{\Et}{Et}
\DeclareMathOperator{\card}{card}
\DeclareMathOperator{\Ens}{Ens}
\DeclareMathOperator{\h}{\underline{H}}
\DeclareMathOperator{\Fonct}{Fonct}
\DeclareMathOperator{\Ker}{Ker}
\DeclareMathOperator{\but}{but}
\DeclareMathOperator{\Ob}{Ob}
\DeclareMathOperator{\ob}{ob}
\DeclareMathOperator{\Lib}{Lib}
\DeclareMathOperator{\esj}{esj}
\DeclareMathOperator{\id}{id}
\DeclareMathOperator{\Id}{Id}
\DeclareMathOperator{\ad}{ad}
\DeclareMathOperator{\Iim}{Im}
\DeclareMathOperator{\pr}{pr}
\DeclareMathOperator{\Cat}{Cat}
\DeclareMathOperator{\car}{car}
\DeclareMathOperator{\SysInd}{SysInd}
\DeclareMathOperator{\Sup}{Sup}
\DeclareMathOperator{\Coker}{Coker}
\DeclareMathOperator{\Filt}{Filt}
\DeclareMathOperator{\acc}{acc}
\DeclareMathOperator{\sur}{sur}
\DeclareMathOperator{\cat}{cat}
\DeclareMathOperator{\Ord}{Ord}
\DeclareMathOperator{\Gr}{Gr}
\DeclareMathOperator{\Cov}{Cov}
\DeclareMathOperator{\Caf}{Caf}
\DeclareMathOperator{\coker}{coker}
\DeclareMathOperator{\coim}{coim}
\DeclareMathOperator{\im}{im}
\DeclareMathOperator{\Ouv}{Ouv}
\DeclareMathOperator{\diag}{diag}
\DeclareMathOperator{\Esp}{Esp}
\DeclareMathOperator{\TOP}{TOP}
\DeclareMathOperator{\Unif}{Unif}
\DeclareMathOperator{\Iind}{\mathbf{Ind}}
\DeclareMathOperator{\Hhomcart}{\mathbf{Homcart}}
\DeclareMathOperator{\Hhomtopess}{\mathbf{Homtopess}}
\DeclareMathOperator{\Ppoints}{\mathbf{Points}}
\DeclareMathOperator{\Ppoint}{\mathbf{Point}}
\DeclareMathOperator{\Ttors}{\mathbf{Tors}}
\DeclareMathOperator{\Ssex}{\mathbf{Sex}}
\DeclareMathOperator{\Ffib}{\mathbf{Fib}}
\DeclareMathOperator{\Vvois}{\mathbf{Vois}}
\DeclareMathOperator{\Ppointess}{\mathbf{Pointess}}
\DeclareMathOperator{\scrHom}{\mathbf{Hom}}
\DeclareMathOperator{\sheafPic}{\mathbf{Pic}}
\DeclareMathOperator{\sheafAut}{\mathbf{Aut}}
\DeclareMathOperator{\Hhomtoplocan}{\mathbf{Homtoplocan}}
\DeclareMathOperator{\Hhomtopan}{\mathbf{Homtopan}}
\DeclareMathOperator{\Pptgeom}{\mathbf{Ptgéom}}
\DeclareMathOperator{\Morsite}{Morsite}
\DeclareMathOperator{\Fib}{Fib}
\DeclareMathOperator{\Fais}{Fais}
\DeclareMathOperator{\kar}{kar}
\DeclareMathOperator{\Groupes}{Groupes}
\DeclareMathOperator{\Restr}{Restr}
\DeclareMathOperator{\Isom}{Isom}
\DeclareMathOperator{\Prol}{Prol}
\DeclareMathOperator{\Res}{Res}
\DeclareMathOperator{\Tors}{Tors}
\DeclareMathOperator{\esp}{esp}
\DeclareMathOperator{\es}{es}
\DeclareMathOperator{\od}{od}
\DeclareMathOperator{\gr}{gr}
\DeclareMathOperator{\sob}{sob}
\DeclareMathOperator{\dfn}{dfn}
\DeclareMathOperator{\Sch}{Sch}
\DeclareMathOperator{\ouv}{ouv}
\DeclareMathOperator{\cons}{cons}
\DeclareMathOperator{\Nerf}{Nerf}
\DeclareMathOperator{\Points}{Points}
\DeclareMathOperator{\Espsob}{Espsob}
\DeclareMathOperator{\Homtop}{Homtop}
\DeclareMathOperator{\Pt}{Pt}
\DeclareMathOperator{\Quot}{Quot}
\DeclareMathOperator{\glob}{glob}
\DeclareMathOperator{\Homfib}{Homfib}
\DeclareMathOperator{\Qouv}{Qouv}
\DeclareMathOperator{\Quotouv}{Quotouv}
\DeclareMathOperator{\ens}{ens}
\DeclareMathOperator{\ab}{ab}
\DeclareMathOperator{\an}{an}
\DeclareMathOperator{\fs}{fs}
\DeclareMathOperator{\red}{red}
\DeclareMathOperator{\Ab}{Ab}
\DeclareMathOperator{\cosq}{cosq}
\DeclareMathOperator{\Ann}{Ann}
\DeclareMathOperator{\Zar}{Zar}


%%%%%%%%%%%%
%%%%% SGA4-2
\DeclareMathOperator{\Sex}{Sex}
\DeclareMathOperator{\vois}{vois}
\DeclareMathOperator{\Vois}{Vois}
\DeclareMathOperator{\Ext}{Ext}
\DeclareMathOperator{\ext}{ext}
\DeclareMathOperator{\Fer}{Fer}
\DeclareMathOperator{\Add}{Add}
\DeclareMathOperator{\Cocart}{Cocart}
\DeclareMathOperator{\cosk}{cosk}
\DeclareMathOperator{\plat}{plat}

\DeclareMathOperator{\faisc}{faisc}
\DeclareMathOperator{\sk}{sk}
\DeclareMathOperator{\ch}{ch}

\def\HH{\mathbf{H}}
\def\qc{\mathrm{qc}}
\def\et{\mathrm{et}}
\def\coh{\mathrm{coh}}

\def\Hvr#1{%
\begin{array}[b]{l}
{\kern2pt\scriptstyle r}\\[-8pt]
{\kern1pt\scriptstyle \vee}\\[-5pt]
H^{#1}
\end{array}}

\def\Hvv#1{%
\begin{array}[b]{l}
\kern2pt\rotatebox{90}{${\scriptscriptstyle\ll}$}\\[-6pt]
H^{#1}
\end{array}}

\def\Hvi#1{%
\begin{array}[b]{l}
{\kern0.2pt\scriptscriptstyle \infty}\\[-8pt]
{\kern1pt\scriptstyle \vee}\\[-5pt]
H^{#1}
\end{array}}

\def\Hv#1{%
\begin{array}[b]{l}
{\kern1pt\scriptstyle \vee}\\[-5pt]
H^{#1}
\end{array}}

\DeclareMathOperator{\cart}{cart}
\DeclareMathOperator{\Ffaisc}{\mathbf{Faisc}}


\def\Ev#1{%
\begin{array}[b]{l}
{\kern1pt\scriptstyle \vee}\\[-5pt]
E'_{#1}
\end{array}}

\DeclareMathOperator{\Cont}{Cont}
\DeclareMathOperator{\Ccont}{\mathscr{C}\!\mathit{ont}}
\DeclareMathOperator{\Site}{Site}
\newcommand\MM{\mathbf{M}}
\def\limtop#1{{\displaystyle{\varprojlim_{#1}\!\mbox{\rm top }}}}
\def\Limtop{\mathop{\varprojlim\!\mbox{\rm top}}}
\DeclareMathOperator{\spec}{spec}
\def\ss{\scriptstyle{\mathscr{S}}}
\DeclareMathOperator{\Gal}{Gal}
\def\smallm{\scriptstyle{\mathscr{M}}}
\DeclareMathOperator{\End}{End}
\DeclareMathOperator{\ass}{ass}
\DeclareMathOperator{\zar}{zar}
\DeclareMathOperator{\Ogg}{Ogg}
\DeclareMathOperator{\res}{res}
\DeclareMathOperator{\tr}{tr}
\DeclareMathOperator{\Supp}{Supp}
\DeclareMathOperator{\codim}{codim}
\DeclareMathOperator{\rk}{rk}
\DeclareMathOperator{\cur}{cur}
\DeclareMathOperator{\cdqc}{cdqc}
\newcommand\DD{\mathbf{D}}
\DeclareMathOperator{\cl}{cl}
\def\JJ{\mathbf{J}}
\def\FF{\mathbf{F}}
\def\BB{\mathbf{B}}
\DeclareMathOperator{\Alget}{Alget}
\DeclareMathOperator{\Ass}{Ass}
\DeclareMathOperator{\supp}{supp}
\DeclareMathOperator{\rad}{rad}
\DeclareMathOperator{\prof}{prof}
\DeclareMathOperator{\cint}{int}
\DeclareMathOperator{\sic}{sic}
\DeclareMathOperator{\Tor}{Tor}
\DeclareMathOperator{\tors}{tors}
\DeclareMathOperator{\Sym}{Sym}
\DeclareMathOperator{\parf}{parf}
\DeclareMathOperator{\tor}{tor}
\DeclareMathOperator{\defn}{defn}
\DeclareMathOperator{\Div}{Div}
\DeclareMathOperator{\Pol}{Pol}
\DeclareMathOperator{\Fex}{Fex}
\DeclareMathOperator{\Lie}{Lie}
\DeclareMathOperator{\tot}{tot}
\DeclareMathOperator{\Aut}{Aut}
\DeclareMathOperator{\HOM}{HOM}
\DeclareMathOperator{\pch}{pch}
\DeclareMathOperator{\Ch}{Ch}
\DeclareMathOperator{\EXT}{EXT}
\DeclareMathOperator{\PIC}{PIC}
\DeclareMathOperator{\Rf}{Rf}
\DeclareMathOperator{\inj}{inj}
\newcommand\VV{\mathbf{V}}

\def\Xv{%
\begin{array}[b]{l}
{\kern1pt\scriptstyle \vee}\\[-5pt]
X
\end{array}}

%%%%%%%
\newcommand\TT{\mathbf{T}}
\newcommand\XX{\mathbf{X}}
\newcommand\EE{\mathbf{E}}
%%%%%%

\def\qedsymbol{\textup{C.Q.F.D.}}
\let\CQFD\qed
\let\cqfd\qed

%%%%%\bigdot
\font\bigsymb=cmsy10 at 4pt
\def\bigdot{{\kern1.2pt\raise 1.5pt\hbox{\bigsymb\char15}}}
%%%%%

\def\Lright{\mathop{\vcenter{\hbox{\hbox{$\longrightarrow$}\lower8pt \hbox{\kern-16pt\hbox{$\longrightarrow$}}}}}}

\def\Aa{{\mathcal A}}
\def\Bb{{\mathcal B}}
\def\Cc{{\mathcal C}}
\def\Dd{{\mathcal D}}
\def\Ee{{\mathcal E}}
\def\Ff{{\mathcal F}}
\def\Gg{{\mathcal G}}
\def\Hh{{\mathcal H}}
\def\Ii{{\mathcal I}}
\def\Jj{{\mathcal J}}
\def\Kk{{\mathcal K}}
\def\Ll{{\mathcal L}}
\def\Mm{{\mathcal M}}
\def\Nn{{\mathcal N}}
\def\Oo{{\mathcal O}}
\def\Pp{{\mathcal P}}
\def\Qq{{\mathcal Q}}
\def\Rr{{\mathcal R}}
\def\Ss{{\mathcal S}}
\def\Tt{{\mathcal T}}
\def\Uu{{\mathcal U}}
\def\Vv{{\mathcal V}}
\def\Ww{{\mathcal W}}
\def\Xx{{\mathcal X}}
\def\Yy{{\mathcal Y}}
\def\Zz{{\mathcal Z}}

\newcommand{\CC}{\mathbf{C}}
\newcommand{\GG}{\mathbf{G}}
\newcommand{\KK}{\mathbf{K}}
\newcommand{\LL}{\mathbf{L}}
\newcommand{\NN}{\mathbf{N}}
\newcommand{\PP}{\mathbf{P}}
\newcommand\QQ{\mathbf{Q}}
\newcommand\RR{\mathbf{R}}
\newcommand\ZZ{\mathbf{Z}}

\def\to{\mathchoice{\longrightarrow}{\rightarrow}{\rightarrow}{\rightarrow}}
\def\from{\mathchoice{\longleftarrow}{\leftarrow}{\leftarrow}{\leftarrow}}
\def\mto{\mathchoice{\longmapsto}{\mapsto}{\mapsto}{\mapsto}}
\def\mfrom{\mathrel{\from\joinrel\mapstochar\,}}
\def\hto{\mathrel{\lhook\joinrel\to}}
\def\hfrom{\mathrel{\from\joinrel\rhook}}

\def\implies{\mathchoice{\Longrightarrow}{\Rightarrow}{\Rightarrow}{\Rightarrow}}
\def\implied{\mathchoice{\Longleftarrow}{\Leftarrow}{\Leftarrow}{\Leftarrow}}
\def\iff{\mathchoice{\Longleftrightarrow}{\Leftrightarrow}{\Leftrightarrow}{\Leftrightarrow}}

\def\To#1{\mathchoice{\xrightarrow{\labelstyle~~#1~}}{\stackrel{#1}{\longrightarrow}}{}{}}
\def\Mto#1{\mathrel{\mapstochar\To{#1}}}
\def\From#1{\mathchoice{\xleftarrow{\labelstyle~~#1~}}{\stackrel{#1}{\longleftarrow}}{}{}}

\newcommand{\isomto}{\overset{\sim}{\longrightarrow}}
\newcommand{\isomfrom}{\overset{\sim}{\longleftarrow}}

\newcommand{\IFF}{\unskip~$\iff$~\ignorespaces}
\newcommand{\THEN}{\unskip~$\implies$~\ignorespaces}

\let\leq\leqslant
\let\geq\geqslant
\let\le\leqslant
\let\ge\geqslant

\def\cf{{cf.}\kern.3em}
\def\Cf{{Cf.}\kern.3em}
\def\ie{{i.e.}\kern.3em}
\def\iev{{i.e.},\xspace}
%\def\loccito{loc.\kern2pt cit.\xspace}
%\def\loccitn{\textit{loc.\kern2pt cit.}\xspace}
%\def\loccit{\sisi{\loccito}{\loccitn}}
\def\EGA{\textup{EGA}\kern.2em}
\def\SGA{\textup{SGA}\kern.3em}

\let\oldS\S
\def\S{\oldS\kern .15em}
\def\ptbl{.\kern .2em}
\def\resp{resp.\kern.3em}
\def\Exp{Exp.\kern.3em}

\let\bar\overline
\let\hat\widehat
\let\tilde\widetilde
\let\dpl\displaystyle
\let\ldots\dots
\let\emptyset\varnothing
\let\skpt\skippointrait
\let\epsilon\varepsilon

\def\sfrac#1#2{{#1}/{#2}}

\DeclareMathOperator{\cd}{cd}
%\DeclareMathOperator{\Hhom}{\mathcal{H}\!\mathit{om}}
\DeclareMathOperator{\Hhom}{\mathbf{Hom}}
%\DeclareMathOperator{\Eext}{\mathcal{E}\!\mathit{xt}}
\DeclareMathOperator{\Eext}{\mathbf{Ext}}
\DeclareMathOperator{\Spec}{Spec}
\DeclareMathOperator{\Tr}{Tr}

\let\emptyset=\varnothing

%%%%%%%%%%%%%%%
%%% expose XVII
%%%%%%%%%%%%%%%

\makeatletter
\let\old@optexpose\optexpose
\newcommand{\empty@optexpose}[1]{}
\newcommand{\paren@optexpose}[1]{\old@optexpose{#1}(}
\newcommand{\refnoexpose}[1]{\let\optexpose\empty@optexpose\ref{#1}\let\optexpose\old@optexpose}%\refnoexpose{X.cor5.2} engendre "5.2" au lieu de "X 5.2"
\newcommand{\neweqref}[1]{\let\optexpose\paren@optexpose\textup{\ref{#1})}\let\optexpose\old@optexpose}%\neweqref{XII.eq4.1} engendre "XII (4.1)" au lieu de "(XII 4.1)"
\makeatother
\newcommand{\refvb}[2][VI${}_{\text{B}}$ ]{\sisi{#1\refnoexpose{#2}}{\ref{#2}}}

\newcommand{\leftexp}[2]{{\vphantom{#2}}^{#1}{#2}}
\newcommand{\AAA}{\mathbf{A}}
\newcommand{\Sf}{\mathfrak{S}}
\DeclareMathOperator*{\quotprojlim}{\text{``$\varprojlim$''}}
\DeclareMathOperator*{\quotinjlim}{\text{``$\varinjlim$''}}
\DeclareMathOperator{\Fl}{F\sisi{\ell}{l}}
\DeclareMathOperator{\Pl}{P\sisi{\ell}{l}}
%\DeclareMathOperator{\cl}{c\sisi{\ell}{l}}
\DeclareMathOperator{\tdf}{tdf}
\DeclareMathOperator{\rg}{rg}
\newcommand{\iInd}{\sisi{I}{i}nd}
\newcommand{\iind}{ind}
\newcommand{\ppro}{pro}
\newcommand{\tTor}{\sisi{T}{t}or}
\newcommand{\ttor}{tor}
\newcommand{\chap}{Chap.}
\newcommand{\Ffex}{\Ff\mathit{ex}}
\newcommand{\sheafk}{\underline{k}}
\newcommand{\sheafn}{\underline{n}}
\newcommand{\sheafGamma}{\underline{\Gamma}}

\sisi{
\DeclareMathOperator{\simarrow}{\rule{.15mm}{1.8mm}\kern -3.4pt\rightsquigarrow}
\newcommand{\sheafell}{\underline{\ell}}
\newcommand{\sheafZZ}{\underline{\ZZ}}
}{
\newcommand{\nderef}[1]{la N.D.E.~(\ref{#1}) page~\pageref{#1}}
}

%%%%%%%%%%%%%%%%%
%%%%% expose XVII
%%%%%%%%%%%%%%%%%



\externaldocument{../01/01.tex}
\externaldocument{../02/02.tex}
\externaldocument{../03/03.tex}
\externaldocument{../04/04.tex}
\externaldocument{../05/05.tex}
\externaldocument{../05bis/05bis.tex}
\externaldocument{../06/06.tex}
%\externaldocument{../07/07.tex}
\externaldocument{../08/08.tex}
\externaldocument{../09/09.tex}
\externaldocument{../10/10.tex}
\externaldocument{../11/11.tex}
\externaldocument{../12/12.tex}
\externaldocument{../13/13.tex}
\externaldocument{../14/14.tex}
\externaldocument{../15/15.tex}
\externaldocument{../16/16.tex}
\externaldocument{../17/17.tex}
\externaldocument{../18/18.tex}
\externaldocument{../19/19.tex}

\setboolean{original}{false}

\begin{document}

\begin{center}
VII. Site et topos étales d'un schéma\\
A. Grothendieck
\end{center}

\selectlanguage{english}version : \showgitstatus\selectlanguage{french}
\tableofcontents

\else
\chapter{Site et topos étales d'un schéma}
%\addtocontents{toc}{A. Grothendieck}
\begin{center}
A. Grothendieck
\end{center}
\fi

\setcounter{pageoriginal}{341}

Dans\pageoriginale le présent exposé et le suivant, nous
développons
les propriétés les plus élémentaires relatives à la
topologie et la cohomologie
étales. Les développements du présent exposé
concernent certaines
propriétés valables pour l'essentiel pour d'autres topologies très
différentes, telle la « topologie fppf ». Dans l'exposé suivant
seront développées des propriétés assez spéciales
à la topologie
étale, tenant à la nature très particulière des morphismes
étales.

Nous suivons partiellement ici trois exposés oraux de J.E. \textsc{Roos} (qui
n'avaient pu être rédigés par lui), notamment dans la démonstration,
de \hyperrefext[VIII][VIII.thm6.3]{6.3}.

Sauf mention expresse du contraire, il sera sous-entendu que les
schémas envisagés dans le présent exposé et les suivants sont éléments
de l'univers fixé $𝒰$.

\section{La topologie étale}\etiquette[1]{VII.sec1}

\subsection{}\etiquette[1.1]{VII.subsec1.1}

Nous désignerons par $(\Sch)$ la catégorie des schémas (éléments de
l'Univers fixé $𝒰$). Rappelons qu'un morphisme $f : X \to Y$ dans
$(\Sch)$ est dit\pageoriginale \emph{étale} s'il est localement de présentation
finie, plat, et si pour tout $y \in Y$, la fibre $X_y$ est discrète,
et ses anneaux locaux sont des extensions finies \emph{séparables} de
$k(y)$. Il revient au même de dire que $f$ est localement de
présentation finie, et que pour tout $Y$ - schéma $Y'$ qui est affine
(au sens absolu) et pour tout sous-schéma $Y'_{0}$ défini par un Idéal
nilpotent, l'application
$$
\Hom_Y (Y', X) \to \Hom_Y (Y'_{0}, X)
$$
est \emph{bijective}. Pour les propriétés les plus importantes de
cette notion, je renvoie à \EGA IV \S\S 17 et 18, et en attendant à
\SGA 1 I et IV.


\subsection{}\etiquette[1.2]{VII.subsec1.2}

On appelle \emph{topologie étale} sur $(\Sch)$ la topologie engendrée
par la prétopologie pour laquelle, pour tout $X \in (\Sch)$,
l'ensemble $\Cov X$ est formé des familles $(X_i \xrightarrow{u_i} X)_{i
\in I}$ (indexées par un $I \in U$), telles que les $u_i$ soient
étales et $X = \bigcup_{i} u_i (X_i)$ (au sens
ensembliste).

Il est commode d'associer à tout $X$ la sous-catégorie $\Et_{/X}$ de
$(\Sch)_{/X}$ formée des flèches $X'\to X$ qui sont étales. On la
munira de la « topologie induite » (III) par la topologie étale de
$(\Sch)$, (appelée encore « topologie étale sur $X$ ») et on
désignera par $X_{\text{ét}}$ le site ainsi obtenu (« site étale » de $X$)
\footnote{Notons cependant qu'il serait préférable de désigner par
  $X_{\text{ét}}$ le topos $\widetilde{X_{\text{ét}}}$ défini par
  le site étale de $X$. Pour des raisons pratiques, nous nous en
  tiendrons dans ce Séminaire aux notations introduites ici (qui sont
  celles du séminaire primitif).}. Notons que tout morphisme de
$\Et_{/X}$ est étale (\SGA 1 I 4.8), d'où s'ensuit qu'une famille
$(X'_i \xrightarrow{u_i} X')_{i \in I}$ dans $X_{\text{ét}}$ est couvrante
si et seulement si elle est \emph{surjective} \ie
$\bigcup_{i} u_{i}(X'_{i}) = X'$. On voit aussitôt
(\cf \hyperref[VII.prop3.1]{3.1}) que $X_{\text{ét}}$ est un $\Uu$-site donc les résultats
des Exp. I à VI sont applicables. Le topos
$\widetilde{X_{\text{ét}}}$ des $𝒰$-faisceaux sur $X_{\text{ét}}$
est le \emph{topos étale} de $X$.

\subsection{}\etiquette[1.3]{VII.subsec1.3}

Pour la\pageoriginale suite, sauf mention expresse du contraire,
toutes les notions
topologiques que nous envisagerons dans $\Et_{/X}$ ou $(\Sch)$,
s'entendent au sens de la topologie étale. D'ailleurs, dans le langage
et les notations, on écrira couramment $X$ au lieu de $X_{\text{ét}}$
ou $Et_{/X}$. Ainsi, on appellera « faisceau sur $X$ » (sous-entendu :
pour la topologie étale) un faisceau sur le site $X_{\text{ét}}$. On
désignera par $\widetilde{X_{\text{ét}}}$ la catégorie de ces
faisceaux, qui est un
$\Uu$-topos. Si $F$ est un faisceau abélien sur $X$, on
désignera par $H^{i} (X, F)$ ses groupes de cohomologie, qui seraient
notés $H^{i} (X_{\text{ét}}, F)$ dans Exp. V.

\subsection{}\etiquette[1.4]{VII.subsec1.4}

Soit $f : X \to Y$ un morphisme de schémas. Alors le foncteur image
inverse
$$
f^{\ast} : Y' \mapsto \fprod{Y'}{X}{Y}
$$
induit un foncteur
\begin{equation*}
f^{\ast}_{\text{ét}} : \Et_{/Y} \to \Et_{/X}\tag{$\ast$}
\end{equation*}
qui commute aux $\varprojlim$ finies et transforme familles couvrantes
en familles couvrantes; c'est par suite un morphisme de sites
$$
f^{\ast}_{\text{ét}} : X_{\text{ét}} \to Y_{\text{ét}}{\quoi},
$$
induisant donc un foncteur sur la catégorie des faisceaux:
$$
f^{\text{ét}}_{\ast} : \widetilde{X_{\text{ét}}} \to \widetilde{Y_{\text{ét}}}
$$
par
$$
f^{\text{ét}}_{\ast}  (F) = F \circ f^{\ast}_{\text{ét}}{\quoi}.
$$
De plus $f^{\text{ét}}_{\ast}$ admet un adjoint à gauche
$$
f^{\ast}_{\text{ét}} : \widetilde{Y_{\text{ét}}} \to
\widetilde{X_{\text{ét}}}
$$
prolongeant\pageoriginale celui envisagé dans $(\ast)$, et commutant aux
$\varinjlim$ quelconques et aux $\varprojlim$ finies \ie
$f^{\ast}_{\text{ét}}$ définit un \emph{morphisme de topos}
$$
f_{\text{ét}} : \widetilde{X_{\text{ét}}} \to \widetilde{Y_{\text{ét}}}{\quoi}.
$$

Évidemment, pour un composé de deux morphisme de schémas
$$
X \xrightarrow{f} Y \xrightarrow{g} Z
$$
on a des isomorphismes canoniques
\begin{equation*}
\begin{cases}
(gf)^{\text{ét}}_{\ast} \simeq g_{\ast}^{\text{ét}}
  f_{\ast}^{\text{ét}} &\quav \mbox{ d'où}\\
(gf)^{\ast}_{\text{ét}} \simeq f^{\ast}_{\text{ét}}
  g^{\ast}_{\text{ét}} &\quav \mbox{ \ie on a
    un isomorphisme}\\
(gf)_{\text{ét}} \simeq g_{\text{ét}} f_{\text{ét}}&\quav
\end{cases}
\end{equation*}
(de sorte qu'on obtient un « pseudo-foncteur » (\SGA 1 VI 8) de la
catégorie $(\Sch)$ dans la catégorie $(\Top)$ des topos $\in 𝒰'$, où
$𝒰'$ est le plus petit univers tel que $𝒰 \in
𝒰'$; comparer IV).

\subsection{Notations}\etiquette[1.5]{VII.subsec1.5}

Dans la suite nous écrirons souvent $f^{\ast}$, $f_{\ast}$ au lieu
de $f^{\ast}_{\text{ét}}$, $f^{\text{ét}}_{\ast}$. Si $f : X \to Y$
est un morphisme de schémas, les foncteurs $R^{i}
f_{\ast}^{\text{ét}}$ seront donc simplement notés $R^{i} f_{\ast}$.
Rappelons avec ces notations la suite spectrale de Leray pour $f$,
et un faisceau abélien $F$ sur $X$ (\hyperrefext[V][V.sec5]{5}):
$$
H^{\ast} (X, F) \Longleftarrow E^{p, q}_2 = H^{p}(Y, R^q f_\ast (F)){\quoi}.
$$
De même si l'on a deux morphismes $f : X \to Y$ et $g : Y \to Z$, on
a une suite spectrale de Leray (suite spectrale de foncteurs
composés)
$$
R^{\ast} (gf)_{\ast} (F) \Longleftarrow E^{p, q}_2 = R^{p} g_{\ast} (R^q
f_\ast (F)){\quoi}.
$$

\subsection{}\etiquette[1.6]{VII.subsec1.6}

Lorsque\pageoriginale $f : X \to Y$ est lui-même un morphisme étale, on peut
considérer $X$ comme objet de $Y_{\text{ét}}$, et on a un \emph{isomorphisme
  de sites} canonique
$$
X_{\text{ét}} \xrightarrow{\sim} (Y_{\text{ét}})_{/X}{\quoi}.
$$
Moyennant cet isomorphisme, le foncteur $f^{\ast}_{\text{ét}} : Y'\mapsto
\fprod{Y'}{X}{Y}$ de $(\ast)$ s'identifie au foncteur changement de base
interne dans la catégorie $\Et_{/Y}$. Il s'ensuit que le foncteur
$$
f^{\ast} : \widetilde{Y_{\text{ét}}} \to \widetilde{X_{\text{ét}}}
$$
est alors isomorphe au foncteur « \emph{restriction} à $Y$ »
$$
f^{\ast} (F) \simeq F \circ i_{X/Y}{\quoi},
$$
où $F$ est un faisceau sur $Y_{\text{ét}}$, et $i_{X/Y} =
X_{\text{ét}} \to Y_{\text{ét}}$
est le foncteur évident (consistant à regarder un schéma étale sur
$X$, $u : X'\to X$, comme un schéma étale sur $Y$ par $fu : X'\to Y$).

\subsection{Questions d'Univers}\etiquette[1.7]{VII.subsec1.7}

On notera que si $X \neq ∅$, alors $X_{\text{ét}}$
\emph{n'est pas} élément de l'univers choisi $𝒰$. Cependant,
comme nous avons signalé, on voit facilement (\hyperref[VII.prop3.1]{3.1}) que
l'on peut trouver une sous-catégorie pleine $C$ de $X_{\text{ét}}$,
\emph{élément de} $𝒰$, satisfaisant aux conditions du « lemme de comparaison » (III), de sorte que le foncteur restriction
induise une équivalence $\widetilde{C} \to
\widetilde{X_{\text{ét}}}$. Ainsi, le \emph{topos étale} est
équivalent à un topos défini en termes d'un site $C \in 𝒰$.
En d'autres termes, $X_{\text{ét}}$ est un $U$-site donc en vertu de
\loccit les résultats des Exposés I à VI sont applicables à ce site.

\section{Exemples de faisceaux}\etiquette[2]{VII.sec2}\pageoriginale

\begin{enumerate}
\renewcommand{\theenumi}{\alph{enumi}}
\renewcommand{\labelenumi}{\theenumi)}
\item Soit $F \in \Ob (\Sch)_{/X}$ un schéma sur $X$, et
  pour tout $X'$ sur $X$, posons
$$
F(X') = \Hom_X (X', F){\quoi}.
$$
Le foncteur $(\Sch)^{\circ}_{/X} \to (\Ens)$ ainsi défini est un
faisceau pour la topologie étale (et même pour la topologie plus
fine fpqc étudiée dans \SGA 3 IV). En d'autres termes la topologie
étale sur $(\Sch)$ est moins fine que la topologie canonique. C'est
en effet, essentiellement, le contenu de \SGA 1 VIII 5.1 (\cf aussi
\SGA 3 IV 6.3.1). A fortiori, la restriction de ce faisceau à
$X_{\text{ét}}$ est un faisceau. On le désignera encore par $F$,
lorsqu'aucune confusion n'est à craindre \footnote{\etiquette[6]{page6}Mais
on fera attention que si $X \neq
  ∅$, le foncteur $\varphi$ qui a $F \in \Ob (\Sch)_{/X}$
  associe le faisceau correspondant $\varphi(F)$ n'est \emph{pas}
  pleinement fidèle, ni même fidèle, et qu'on ne peut reconstituer $F$
(mod. isom.) connaissant $\varphi(F)$. Par exemple; si $S = \Spec k$,
  $k$ corps alg. clos, la connaissance de $\varphi(F)$ équivaut à
  celle de l'ensemble sous-jacent à $F$ seulement (en vertu de
  \hyperrefext[VIII][VIII.cor2.4]{2.4}).
Comparer IV}. Noter aussi que le foncteur ainsi obtenu
\begin{equation*}
(\Sch)_{/X} \to \widetilde{X_{\text{ét}}}\tag{$\ast$}
\end{equation*}
commute aux $\varprojlim$ finies (c'est trivial). Cela implique par
exemple que lorsque $F$ est un schéma en groupes (\resp\ldots) sur $X$,
alors le faisceau qu'il définit est un faisceau en groupes (\resp\ldots).
 Notons que le foncteur $\Et_{/X} \to \widetilde{X_{\text{ét}}}$
induit par (*) n'est autre que le foncteur canonique, associant à
tout $X'\in \Ob \widetilde{X_{\text{ét}}}$ le foncteur sur
$X_{\text{ét}}$ qu'il représente. C'est donc un isomorphisme de la
catégorie $\Et_{/X}$ sur une sous-catégorie pleine du topos étale
$\widetilde{X_{\text{ét}}}$, par laquelle nous identifions
généralement un $X \in \Ob X_{\text{ét}}$  au faisceau
correspondant, qui sera noté $\widetilde{X}$ ou simplement $X$.

Évidemment on aura $H^{\circ}(X, F) = \Hom_{(\Sch)_{/X}} (X, F)$.

Donnons également une interprétation de $H^1(X, F)$ lorsque $F$ est
un préschéma en groupes sur $X$ (commutatif si l'on veut, de sorte
que la définition\pageoriginale de $H^{1}(X, F)$ relève de Exp. V;
pour le cas général on pourra consulter la thèse de J. Giraud
\footnote{\etiquette[7]{page7}}). Alors des raisonnements bien connus, que
je me dispense de répéter ici, montrent que $H^1 (X, F)$ est
canoniquement isomorphe au groupe des classes (mod. isomorphisme) de
\emph{faisceaux d'ensembles sur} $X_{\text{ét}}$
  \emph{principaux homogènes sous $F$\footnote{ou $F$-torseurs dans la
    terminologie de \loccit.}}. Lorsque $F$ est \emph{affine} sur
$X$, alors \SGA 1 VIII 2.1 implique (toujours par des arguments
standards, \cf thèse de J. \textsc{Giraud}) que $H^{1} (X, F)$ est aussi le
groupe des classes de \emph{schémas} $P$ sur $X$, sur lesquels $F$
opère à droite et qui sont « fibrés principaux homogènes sur $X$ au
sens de la topologie étale » \ie localement triviaux dans le sens de
la topologie étale \footnote{ou $F$-\emph{torseurs représentables}
  dans la terminologie maintenant reçue}.
 \end{enumerate}

\begin{enonce*}[remark]{Remarques 2.1}\etiquette[2.1]{VII.rem2.1}
On fera attention que si l'on désigne, pour tout schéma $Z$ sur $X$,
par $\varphi_{X} (Z)$ le faisceau associé sur $X_{\text{ét}}$, et si $f : Y
\to X$ est un morphisme de schéma, on a un homomorphisme évident
(fonctoriel en $Z$)
\begin{equation*}
f^{\ast} (\varphi_X (Z)) \to \varphi_Y (Z_Y),\quad Z_Y =
\fprod{Z}{Y}{X}{\quoi},\tag{$\ast$}
\end{equation*}
\ie un homomorphisme évident
$$
\varphi_X (Z) \to f_{\ast} (\varphi_Y (Z_Y)){\quoi},
$$
savoir l'homomorphisme fonctoriel en $X'\in \Ob X_{\text{ét}}$
$$
\Hom_X (X', Z) \to \Hom_Y (X'_{Y}, Z_{Y}){\quoi}.
$$
Mais on fera attention qu'\emph{en général $(\ast)$ n'est pas un
isomorphisme, \ie la formation du faisceau étale associé à un schéma
relatif ne commute pas aux foncteurs images inverses}. Cependant
$(\ast)$ est un isomorphisme dans le cas particulier où $Z$ est étale
sur $X$. De même,\pageoriginale le foncteur
$$
\varphi_{X} : \Sch_{/X} \to \widetilde{X_{\text{ét}}}
$$
\emph{n'est pas fidèle} si $X \neq ∅$ \footnote{\Cf pour ceci
  la note au bas de la page \pagehyperref[page6]{6}}; cependant sa restriction à $X_{\text{ét}}$,
qui n'est autre que le foncteur canonique
$X_{\text{ét}}\to \widetilde{X_{\text{ét}}}$, est \emph{pleinement
  fidèle}, car en vertu de a) la topologie de $X_{\text{ét}}$ est
moins fine que sa topologie canonique.

\begin{enumerate}
\renewcommand{\theenumi}{\alph{enumi}}
\renewcommand{\labelenumi}{\theenumi)}
\setcounter{enumi}{1}
\item Notons que $(\ast)$ commute également aux sommes directes
  indexées par les $I \in 𝒰$, comme on vérifie facilement. En
  particulier, si pour tout ensemble $I$, on désigne par $I_{X}$ le $X$
  schéma constant défini par $I$ (\SGA 3 I 1.8), somme directe de I
  copies de $X$, alors le faisceau associé n'est autre que le faisceau
  constant $I_{X_{\text{ét}}}$ (somme directe de $I$ copies du  faisceau
  final sur $X_{\text{ét}}$). Comme le foncteur $I \mapsto I_X$ commute
  également aux $\varprojlim$ finies, on voit qu'il transforme groupe
  en groupe, groupe commutatif en groupe commutatif \etc Si $G$ est
  un groupe commutatif ordinaire, on écrira simplement $H^{i}(X, G)$
  au lieu de $H^{i}(X, G_X)$. Supposons, par exemple, que $G$ soit un
  groupe fini, alors $G_X$ est fini donc affine sur $X$, donc en
  utilisant la remarque finale de a) on obtient une interprétation de
  $H^{1} (X, G)$ comme l'ensemble \emph{des classes de revêtements
  principaux galoisiens de groupe $G$} (\SGA 1 V 2.7). Lorsque $X$ est
  connexe et muni d'un point géométrique $a$, alors en termes du «   groupe fondamental » $\pi_1 (X, a)$ (\SGA 1 V 7) on obtient
  l'isomorphisme canonique
\begin{equation*}
H^1(X, G) \simeq \Hom(\pi_1(X, a), G){\quoi},\tag{${\ast}{\ast}$}
\end{equation*}
(où on a supposé $G$ commutatif).

On peut dire qu'en passant de la cohomologie de Zariski à la topologie
étale, « on a fait ce qu'il fallait » pour obtenir « le bon »
$H^1$ (qui figure au 2ème\pageoriginale membre de $(\ast)$) pour un groupe de
coefficients constant fini $G$. \emph{C'est un fait remarquable, qui sera
démontré par la suite de ce séminaire, que cela suffit également pour
trouver les « bons » $H^{i} (X, G)$  pour tout groupe de coefficients
de torsion (du moins si $G$ est premier aux caractéristiques
résiduelles de $X$)}.

\item Soit $F$ un $\Oo_X$-Module sur $X$, au sens de la topologie de
Zariski. Alors (avec les notations de \SGA 3 I 4.6) on définit un
foncteur
$$
W(F) : (\Sch)_{/X}^\circ \to (\Ens)
$$
par
$$
W(F)(X') = \Gamma (X', F \otimes_{\Oo_{X}} \Oo_{X'}){\quoi}.
$$
Ce foncteur est encore un faisceau pour la topologie étale (et même
pour la topologie fpqc) comme il résulte encore de \SGA 1 VIII 1.6 et
\SGA 3 IV 6.3.1. A fortiori, la restriction de $W(F)$ à $X_{\text{ét}}$ est
encore un faisceau, qu'on notera encore $W(F)$. Par définition on aura
donc $H^{\circ} (X, W(F)) = \Gamma(F) = H^{\circ} (X, F)$. Mais on a
mieux si $F$ est quasi-cohérent, \cf \hyperref[VII.prop4.3]{4.3}
\end{enumerate}
\end{enonce*}

\section{Générateurs du topos étale. Cohomologie d'une
  $\protect\varinjlim$ de
  faisceaux}\etiquette[3]{VII.sec3}

\begin{enonce*}{Proposition 3.1}\etiquette[3.1]{VII.prop3.1}
Soient $X$ un schéma, $C$ une sous-catégorie pleine de $X_{\text{\rm
    ét}}$, telle
que pour tout $X'$ étale sur $X$ qui est affine, $C$ contienne un
objet isomorphe à $X'$. Alors $C$ est une « famille de générateurs
topologiques » du site $X_{\text{\rm ét}}$ donc une famille de générateurs du
topos $\widetilde{X_{\text{\rm ét}}}$, et le foncteur restriction
$\widetilde{X_{\text{\rm ét}}} \to \widetilde{C}$  est une équivalence de
catégories (où $C$ est muni de la topologie induite par celle de
    $X_{\text{\rm ét}}$).
\end{enonce*}

Trivial\pageoriginale à l'aide du « lemme de comparaison ».

\begin{enonce*}{Corollaire 3.2}\etiquette[3.2]{VII.cor3.2}
Supposons que $X$ soit quasi séparé. Appelons site étale restreint
de $X$ la sous-catégorie pleine $C$ de $X_{\text{\rm ét}}$ formée
des schémas étales sur $X$ qui sont de présentation finie sur $X$,
munie de la topologie induite par $X_{\text{\rm ét}}$. Alors :
\begin{enumerate}
\renewcommand{\theenumi}{\roman{enumi}}
\renewcommand{\labelenumi}{(\theenumi)}
\item $C$ est stable par produits fibrés, et est un site de type fini
si $X$ est quasi-compact.

\item Le foncteur restriction $\widetilde{X_{\text{\rm ét}}} \to
\widetilde{C}$ est une équivalence de catégories.
\end{enumerate}
\end{enonce*}

L'assertion (i) est triviale, et (ii) résulte de \hyperref[VII.prop3.1]{3.1} car $C$ satisfait
à la condition de \hyperref[VII.prop3.1]{3.1} grâce au fait que $X$ est quasi-séparé, qui
implique que si $X'$ est quasi-compact, par exemple affine, il est
quasi-compact \emph{sur} $X$ (\EGA IV 1.2.4, où on fait $Z = \Spec
\ZZ$), donc de présentation finie sur $X$ si $X'$ est
localement de type fini sur $X$.

\begin{enonce*}{Proposition 3.3}\etiquette[3.3]{VII.prop3.3}
Supposons $X$ quasi-compact et quasi-séparé. Alors les foncteurs
$H^{q} (X_{\text{\rm ét}}, F)$ sur la catégorie des faisceaux
abéliens sur $X_{\text{\rm ét}}$ commutent aux $\varinjlim$.
\end{enonce*}

En effet, on peut remplacer $X_{\text{ét}}$ par $C$ en vertu de
\hyperref[VII.cor3.2]{3.2} (ii), or comme $X$ est quasi-compact on a $X \in
\Ob C$, et la conclusion résulte de \hyperref[VII.cor3.2]{3.2} (i) et
\hyperrefext[VI][VI.subsubsec6.1.2]{6.1.2} (3) .

\section{Comparaison avec d'autres topologies}\etiquette[4]{VII.sec4}

\setcounter{subsection}{-1}
\subsection{}\etiquette[4.0]{VII.subsec4.0}

Tout d'abord notons que les exemples de faisceaux sur
$X_{\text{ét}}$ considérés au $N^{\circ}$ 2 sont en fait de façon
naturelle des \emph{restrictions} de faisceaux\pageoriginale définis
sur $(\Sch)_{/X}$ muni de sa topologie étale (ou même de la
topologie fpqc). Soit de façon générale
$$
u^{\ast} : X_{\text{ét}} \to (\Sch)_{/X}
$$
le foncteur d'inclusion, qui est continu (III) et commute aux
$\varprojlim$ finies, donc définit un \emph{morphisme de sites}
$$
u : (\Sch)_{/X} \to X_{\text{ét}}{\quoi},
$$
d'où un foncteur
$$
u_{\ast} : \widetilde{(\Sch)_{/X}} \to \widetilde{X_{\text{ét}}}
$$
sur les catégories de faisceaux associées:
$$
u_{\ast} (F) \to F \circ u{\quoi},
$$
et un foncteur adjoint à gauche\footnote{En toute rigueur, comme
$(\Sch)_{/X}$ est une $\Uu$-catégorie mais \emph{pas} un
$\Uu$-site on ne peut invoquer III pour l'existence de
$u^{\ast}$; Cependant, on construit facilement $u^{\ast}$ par
$$
u^{\ast} (F) (X') = (p_{X'})^{\ast}_{\text{ét}} (F) (X')\quad (X'\in \Ob \Sch_{/X}){\quoi},
$$
où $p_{X'}$ désigne le morphisme structural $X'\to X$. D'autre
part, pour donner un sens à \hyperref[VII.prop4.1]{4.1} (ii) et justifier la
démonstration indiquée de \hyperref[VII.prop4.1]{4.1}, il y a lieu
d'introduire un univers $\Vv$ tel que $\Uu \in \Vv$, et de
considérer les faisceaux intervenant dans les formules
\hyperref[VII.prop4.1]{4.1} (i) comme des faisceaux à valeurs dans
$\Vv\text{-}(\Ens)$. On peut aussi, si on ne veut travailler qu'avec
des $U$-sites, remplacer $(\Sch)_{/X}$ par une sous-catégorie pleine
$(\Sch)_{/X}$ stable par $\varprojlim$ finies, contenant $C$ de
\hyperref[VII.prop3.1]{3.1}, et qui soit $U$-petite.} de ce dernier
$$
u^{\ast} : \widetilde{X_{\text{ét}}} \to \widetilde{\Sch_{/X}}
$$
Ceci posé:

\begin{enonce*}{Proposition 4.1}\etiquette[4.1]{VII.prop4.1}
\begin{enumerate}
\renewcommand{\theenumi}{\roman{enumi}}
\renewcommand{\labelenumi}{(\theenumi)}
\item Le foncteur $u^{\ast}$ est pleinement fidèle, donc pour tout
faisceau $G$ sur $X_{\text{\rm ét}}$, l'homomorphisme canonique $G
\mapsto u_{\ast} u^{\ast} (G)$ est un isomorphisme.

\item Soient $G$ (\resp $F$) un faisceau abélien sur $X_{\text{\rm ét}}$ (\resp
$(\Sch)_{/X}$). Alors les homomorphismes canoniques (définis par
exemple comme edge-homomorphismes de la suite spectrale de Leray
suivants sont\pageoriginale des isomorphismes:
\begin{gather*}
H^{\ast} (X_{\text{\rm ét}}, F \circ u) \xrightarrow{\sim} H^{\ast}((\Sch)_{/X}, F)\\
H^{\ast} (X_{\text{\rm ét}}, F \circ u) \xrightarrow{\sim} H^{\ast}((\Sch)_{/X}, F){\quoi}.
\end{gather*}
\end{enumerate}

\end{enonce*}

En effet, le foncteur $u$ est un comorphisme et on peut alors
appliquer (III).

Ainsi, (i) montre que pour l'étude de la cohomologie des faisceaux,
il est essentiellement équivalent de travailler avec le
« petit » site étale $X_{\text{ét}}$, ou le « gros » site étale
$(\Sch)_{/X}$.

\setcounter{subsection}{1}
\subsection{}\etiquette[4.2]{VII.subsec4.2}

D'autre part, il y a lieu d'introduire dans $(\Sch)$ (donc dans les
$(\Sch)_{/X}$) diverses autres topologies que la topologie étale, dont
les plus utiles sont définies dans \SGA 3 IV 6.3. La plus grossière
parmi ces dernières est la \emph{topologie de Zariski} $(\Zar)$, définie
par la prétopologie où les familles couvrantes sont les familles
surjectives d'immersions ouvertes; elle est moins fine que la
topologie étale. La plus fine de ces topologies est la \emph{topologie « fidèlement plat quasi-compacte »}, en abrégé (fpqc), qui est la moins
fine des topologies pour lesquelles les familles couvrantes au sens de
Zariski, ainsi que les morphismes fidèlement plats quasi-compacts,
sont couvrants; la topologie fpqc est plus fine que la topologie
étale. Comme nous l'avons déjà remarqué, les divers exemples de
faisceaux sur $(\Sch)/X$ envisagés au $N^{\circ}$ 2 sont en fait déjà
des faisceaux pour la topologie fpqc.

\subsubsection{}\etiquette[4.2.1]{VII.subsubsec4.2.1}

On fera attention cependant que pour un faisceau abélien $F$ sur
$(\Sch)_{/X}$ pour la topologie (fpqc), (ou une topologie, telle
fppf, envisagée dans \SGA 3, les  groupes de cohomologie de $F$ pour
la topologie fpqc (\resp\ldots) ne sont pas toujours isomorphes aux
groupes de cohomologie pour\pageoriginale la topologie étale, et
ceci même si $X$ est le spectre d'un corps $k$, et $F$ est
représentable par un groupe algébrique sur $k$, même en ce qui
concerne le $H^{1}$. De façon générale, on peut montrer que la
topologie étale donne les « bons » groupes de cohomologie pour
les groupes de coefficients qui sont des schémas en groupes
\emph{étales}, ou plus généralement \emph{lisses}, sur
$X$\footnote{\Cf A Grothendieck, Le groupe de Brauer III th 11.7, in
« Dix exposés sur la topologie des schémas », North Holland Pub.
Cie}, mais il n'en est plus de même pour des schémas en groupes tels
que les groupes \emph{radiciels} sur $S$, pour lesquels il y a lieu
de remplacer la topologie étale, encore trop grossière, par la
topologie fpqc, ou fppf.

\subsubsection{}\etiquette[4.2.2]{VII.subsubsec4.2.2}

Comme exemple des relations entre les cohomologies relatives à des
topologies différentes, signalons ici le cas des topologies de Zariski
et de la topologie étale. Nous désignerons par $X_{\Zar}$ le site des
ouverts de Zariski de $X$, de sorte qu'on a un foncteur d'inclusion
canonique
$$
u^{\ast} : X_{\Zar} \to X_{\text{ét}}
$$
qui définit un morphisme de sites
$$
u : X_{\text{ét}} \to X_{\Zar}{\quoi},
$$
d'où des foncteurs correspondants image directe
$$
f_{\ast} : \widetilde{X_{\text{ét}}} \to \widetilde{X_{\Zar}}\quad
(f_{\ast} (F) = F \circ u){\quoi},
$$
et image inverse
$$
f^{\ast} : \widetilde{X_{\Zar}} \to \widetilde{X_{\text{ét}}}{\quoi},
$$
adjoints l'un de l'autre; géométriquement, il y a lieu d'interpréter
le couple $(f_{\ast}, f^{\ast})$ comme un morphisme de topos
$$
f : \widetilde{X_{\text{ét}}} \to \widetilde{X_{Zar}}{\quoi}.
$$
On en déduit un homomorphisme de foncteurs cohomologiques
$$
H^{\ast} (X_{\Zar}, f_{\ast} (F)) \to H^{\ast} (X_{\text{ét}}, F)
$$
et\pageoriginale une suite spectrale de Leray
$$
H^{\ast} (X_{\text{ét}}, F) \Longleftarrow H^{p} (X_{\Zar}, R^q
f_{\ast} (F)){\quoi},
$$
où $F$ est un faisceau abélien sur $X_{\text{ét}}$. Cette suite spectrale
résume les relations générales entre cohomologie étale et cohomologie
de Zariski. Bien entendu, pour des faisceaux de coefficients $F$ tel
que des faisceaux de coefficients constants, cette suite spectrale en
général est loin d'être triviale, \ie en général on aura $R^q f_{\ast}
(F) \neq 0$  (prendre notamment le cas où $X$ est le spectre d'un
corps). Cependant:

\begin{enonce*}{Proposition 4.3}\etiquette[4.3]{VII.prop4.3}
Soit $F$ un faisceau de modules variable sur le préschéma $X$
(faisceau au sens de la topologie de Zariski), d'où un faisceau
$F_{\text{\rm ét}}$ sur $X_{\text{\rm ét}}$ (\cf 2, c) et un
homomorphisme de foncteurs cohomologiques
$$
H^{\ast} (X, F) \to H^{\ast} (X_{\text{\rm ét}}, F_{\text{\rm ét}}){\quoi}.
$$
Lorsque $F$ est quasi-cohérent, l'homomorphisme précédent est un isomorphisme.
\end{enonce*}

En effet, le premier membre n'est autre que $H^{\ast} (X_{\Zar},
f_{\ast} (F_{\text{ét}}))$, et avec les notations précédentes, il
suffit de prouver qu'on a
$$
R^q f_{\ast} (F) = 0\quad \mbox{pour}\quad q > 0{\quoi}.
$$
Cela résulte aussitôt, grâce au procédé de calcul des $R^q f_{\ast}$
et grâce au fait que les ouverts affines forment une base de la
topologie de $X$, du

\begin{enonce*}{Corollaire 4.4}\etiquette[4.4]{VII.cor4.4}
Si $X$ est affine, on a
$$
H^{q} (X_{\text{\rm ét}}, F_{\text{\rm ét}}) = 0\quad \mbox{pour}\quad q > 0{\quoi}.
$$
\end{enonce*}

Pour le voir, soient $C = X_{\text{ét}}$, $C'$ la sous-catégorie pleine
formée des $X'$ \emph{affines}, alors en vertu de \hyperref[VII.prop3.1]{3.1} on aura
$$
H^{q} (C, F_{\text{ét}}) \simeq H^{q} (C', F | C'){\quoi}.
$$

Il suffit\pageoriginale donc de prouver que $F | C'$ est
$C'$-acyclique \hyperrefext[V][V.def4.1]{4.1}, ou encore satisfait la condition de
\hyperrefext[V][V.prop4.3]{4.3} \ie que pour tout $X'\in \Ob C'$ et toute famille
couvrante $ℛ = (X'_{i} \to X')_i$ dans $C'$, on a $H^{q}
(ℛ, F) = 0$ pour $q > 0$. Or on peut supposer $(X'_{i})_{i}$
fini, puis, quitte à remplacer les $X'_{i}$ par leur somme, que la
famille couvrante consiste en un morphisme $X'_{i} \to X'$ qui est
couvrant, \ie (étale et) surjectif. Donc on est ramené à prouver que
si $f : X \to Y$ est un morphisme étale surjectif de schéma affines
et $F$ un module quasi-cohérent sur $X$, alors
\begin{equation*}
H^q (X /Y, F) = 0\quad \mbox{pour}\quad q > 0{\quoi}.\tag{$\ast$}
\end{equation*}
Or ceci est démontré dans TDTE I, B, \hyperref[VII.subsec1.1]{1.1}. (dans FGA, \cf réf [4] de
IX).

\begin{enonce*}[remark]{Remarques 4.5}\etiquette[4.5]{VII.rem4.5}
En fait, par \loccit, pour avoir $(\ast)$ il suffit que $X \to Y$ soit
surjectif, et plat (au lieu de étale). Cela permet de montrer par la
démonstration  précédente, qu'on a encore des isomorphismes analogues
à celui de \hyperref[VII.prop3.3]{3.3}, en y remplaçant le site étale $X_{\text{ét}}$ par
$(\Sch)_{/X'}$ muni d'une quelconque des topologies plus fines que celle
de Zariski envisagées dans SGA 3 IV 6.3, par exemple la topologie fpqc.
\end{enonce*}

\section{Cohomologie d'une limite projective de schémas}\etiquette[5]{VII.sec5}

\subsection{}\etiquette[5.1]{VII.subsec5.1}

Soit $I$ un ensemble préordonné filtrant croissant,
$$
\mathfrak{X} = (X)_{i \in I}
$$
un système projectif de schémas, les morphismes de transition $u_i :
X_i \to X_j$ $(i \geq j)$ étant \emph{affines}. Rappelons (\EGA IV
8\ldots) que sous ces conditions,\pageoriginale la limite projective
$$
X  = \varprojlim X_i
$$
existe dans la catégorie des schémas, (et même dans la catégorie des
espaces annelés en anneaux locaux) et peut se construire ainsi:  on
choisit $i_{0} \in I$, de sorte que pour $i \geq i_0$, on a un
$X_{i_0}$ - isomorphisme (\EGA II 1.2 et 1.3)
$$
X_i = \Spec 𝒜_{i}{\quoi},
$$
où $𝒜_i$ est une Algèbre quasi-cohérente sur $X_i$, les
$𝒜_i$ $(i \geq i_0)$ formant donc un système inductif de
telles Algèbres. Posant
$$
𝒜 = \varinjlim_{i} 𝒜_i{\quoi},
$$
on peut prendre
$$
X = \Spec 𝒜{\quoi}.
$$
D'ailleurs, désignant par esp $S$ l'espace topologique sous-jacent à
un schéma $S$, on montre (\EGA IV 8\ldots) que l'application
canonique
$$
\esp X \to \varprojlim \esp (X_i)
$$
est un \emph{homéomorphisme} et que l'homomorphisme canonique de
faisceaux d'anneaux
$$
\varinjlim_{i} \gamma^{-1}_{i} (\Oo_{X_i}) \to \Oo_X
$$
est bijectif, où $\gamma_i : X \to X_i$ est l'application continue
canonique.

\subsection{}\etiquette[5.2]{VII.subsec5.2}

De façon imagée, on peut résumer le contenu de \EGA IV 8, 9 en
disant que, \emph{lorsque $X_{i_0}$ est quasi-compact et
quasi-séparé}, alors « toute donnée de nature schématique sur $X$,
de présentation finie sur $X$ », est équivalente\pageoriginale à une
donnée de même
nature sur un des $X_i$, « pour $i$ assez grand ». Ainsi, on prouve
(\EGA IV 8 $\ldots$):

\begin{enumerate}
\renewcommand{\theenumi}{\alph{enumi}}
\renewcommand{\labelenumi}{(\theenumi)}
\item si $i_1 \in I$ et si $X'_{i_1}$, $X''_{i_1}$ sont deux
schémas de présentation finie sur $X_{i_1}$, alors posant pour tout $i
\geq i_1$
$$
X'_{i} = \fprod{X'_{i_1}}{X_i}{X_{i_1}},\quad X' = \fprod{X'_{i_1}}{X}{X_{i_1}}{\quoi},
$$
et définissant de même $X''_i$, $X''$, l'application canonique
$$
\varinjlim_{i \geq i_1} \Hom_{X_i} (X'_i, X''_i) \to
\Hom_X (X', X'')
$$
est bijective.

\item Pour tout schéma $X'$ de présentation finie sur $X$, il existe
un indice $i_1 \in I$, un schéma $X'_{i_1}$ de présentation finie
sur $X_{i_1}$, et un $X$-isomorphisme
$$
X'\simeq \fprod{X'_{i_1}}{X}{X_{i_1}}{\quoi}.
$$
\end{enumerate}

\subsection{}\etiquette[5.3]{VII.subsec5.3}

On peut exprimer le résultat précédent, et les résultats de nature
analogue contenus dans EGA IV 8, dans le langage des $\varinjlim$ de
catégories fibrées introduit dans VI. Pour ceci, considérons la
catégorie $\Ff$ des morphismes $f : T \to S$ de présentation
finie de schémas, et le « foncteur but »
$$
\pi : \Ff \to (\Sch){\quoi},
$$
qui est évidemment un \emph{foncteur fibrant}, les catégories fibres
étant d'ailleurs équivalentes à des catégories
$\Uu$-petites. Considérons l'image inverse de la catégorie
fibrée $\Ff/(\Sch)$ par le foncteur
$$
i \mapsto X_i : \cat (I) \to (\Sch)
$$
(où\pageoriginale $\cat(I)$ est la catégorie associée à $I$, avec
$\Hom(i, j) \neq ∅$ si et seulement si $i \geq j$), d'où une
catégorie fibrée
$$
\pi_{\mathfrak{X}} : \Ff_{\mathfrak{X}} \to \cat(I){\quoi},
$$
dont la fibre en chaque $i \in I$ est canoniquement isomorphe à la
sous-catégorie pleine $\Ff_{X_i}$ de $(\Sch)_{/X_i}$ formée des
schémas de présentation finie sur $X_i$, le morphisme de changement de
base $\Ff_i \to \Ff_{j}$ relatif à $j \geq i$ n'étant
autre que $X'_i \mapsto X'_j = \fprod{X'_i}{X_j}{X_i}$. Ceci
posé, associant à tout $X'_i$ de présentation finie sur $X_i$ son
image inverse sur $X$
$$
\varphi(X'_i) = \fprod{X'_i}{X}{X_i}{\quoi},
$$
on trouve un foncteur naturel
$$
\varphi : \Ff_{\mathfrak{X}} \to \Ff_{X'}{\quoi},
$$
qui transforme évidemment morphisme cartésien de
$\Ff_{\mathfrak{X}}$ en isomorphisme, donc par définition se
factorise de façon unique par un foncteur canonique
\begin{equation*}
\psi : \varinjlim_{\Ff_X /\Cat I}
\Ff_{\mathfrak{X}} \to \Ff_{X}{\quoi}.\tag{$\ast$}
\end{equation*}
Compte tenu de la description du premier membre, le résultat de \EGA IV
8 rappelé plus haut peut s'énoncer alors en disant que $(\ast)$ est
une \emph{équivalence de catégories}.

\subsection{}\etiquette[5.4]{VII.subsec5.4}

L'énoncé a) de \hyperref[VII.subsec5.2]{5.2}. ci-dessus peut se préciser de diverses façons,
en introduisant quelque ensemble $(\Mm)$ de morphismes de
schémas, stable par changement de base, et en énonçant que pour un
$X_{i_1}$-morphisme donné\pageoriginale $u_{i_1} : X'_{i_1} \to
X''_{i_1}$, le
morphisme correspondant $u': X'\to X''$ est $\in
(\Mm)$ si et seulement si il existe $i\geq i_{1}$ tel que le
$X_{i}$-morphisme $u_{i}:X'_{i}\to X''_{i}$ soit $\in
(\Mm)$. L'énoncé obtenu ainsi est vrai par exemple lorsque
$(\Mm)$ est l'un des ensembles de flèches suivants: morphismes
propres (respectivement projectifs, \resp quasi projectifs, \resp affines, \resp quasi affines, \resp finis, \resp quasi finis, \resp
radiciels, \resp surjectifs, \resp plats, \resp lisses, \resp étales,
\resp non ramifiés). Le lecteur trouvera les énoncés correspondants
dans \EGA IV par 8 (pour les 9 premiers) par. 11 (pour le cas plat) par.
17 (pour les trois derniers cas).

\subsection{}\etiquette[5.5]{VII.subsec5.5}

Remplaçons alors $\Ff \to (\Sch)$ par la sous catégorie
fibrée $\Gg$ formée des \emph{morphismes étales} de
présentation finie, dont la catégorie fibre $\Gg_{S}$, pour
tout schéma $S$, n'est autre que la catégorie des schémas étales de
présentation finie sur $S$, et formons de même la catégorie fibrée
$\Gg_{\mathfrak{X}}$ sur $\cat(I)$, dont la catégorie fibre en
tout $i \in I$ est la catégorie $\Gg_{X_i}$. En vertu de \hyperref[VII.cor3.2]{3.2}
(ii), les topos étales $\widetilde{X_{i_{\text{ét}}}}$,
$\widetilde{X_{\text{ét}}}$ sont aussi canoniquement équivalents à
$\widetilde{\Gg_{X_{i}}}$,
$\widetilde{\Gg_{X}}$, (où tout $\Gg_S$ est muni
de la topologie induite par $S_{\text{ét}}$). De plus, par \hyperref[VII.cor3.2]{3.2} (i), chaque
$\Gg_{X_i}$ est un site de type fini, d'ailleurs équivalent à
un site $\Uu$-petit. Compte tenu des topologies sur les
$\Gg_{X_i}$, $\Gg$ devient un \emph{site fibré} sur
$\cat(I)$. Ceci posé, on peut énoncer le

\begin{enonce*}{Lemme 5.6}\etiquette[5.6]{VII.lem5.6}
Le foncteur canonique
$$
\varinjlim_{\Gg_{\mathfrak{X}}/\cat(I)}
\Gg_{\mathfrak{X}} \to \Gg_X
$$
définit\pageoriginale une équivalence, respectant les topologies, du site
$\varinjlim$ des sites étales restreints (\hyperref[VII.cor3.2]{3.2}) des $X_i$, avec le site
étale restreint de $X = \varprojlim_i X_i$.
\end{enonce*}

Le fait qu'on obtient une équivalence de catégories est l'un des
énoncés rappelés dans \hyperref[VII.subsec5.4]{5.4} (celui où $(\Mm)$ est
l'ensemble des morphismes étales dans $(\Sch)$). L'assertion
relative aux topologies s'obtient de même en prenant $(\Mm)=$
ensemble des morphismes surjectifs dans $(\Sch)$.

Le résultat précédent nous permet d'appliquer les résultats de à la
situation présente. On trouve en particulier

\begin{enonce*}{Théorème 5.7}\etiquette[5.7]{VII.thm5.7}
Soit $\mathfrak{X} = (X_i)_{i \in I}$ un système projectif de
schémas, avec $I$ ensemble préordonné filtrant croissant. Supposons
que les morphismes de transition $u_{ij} : X_j \to X_i$ sont affines
(de sorte que le schéma $X = \varprojlim X_i$ est défini
(\hyperref[VII.subsec5.1]{5.1})), et que les $X_i$ sont quasi-compacts et
quasi-séparés. Considérons le site fibré $\Gg$ sur $\cat(I)$
explicité dans \hyperref[VII.subsec5.5]{5.5}, dont la fibre $\Gg_i$ en $i \in
I$ est canoniquement isomorphe au « site étale restreint » de
$X_i$ (formé des schémas étales de présentation finie sur $X_i$).
Soit $F$ un faisceau abélien sur $\Gg$ (\ie $(3)$) un foncteur
$\Gg^{0} \to (\Ab)$ dont la restriction $F_i$ à chaque $\Gg_i$ est
un faisceau, \ie un faisceau sur $X_i$). Soit $F_{\infty} =
\varinjlim_{i} u^{\ast}_i (F_i)$ le faisceau induit sur $X$, où
$u_i : X \to X_i$ est le morphisme canonique. Alors les
homomorphismes canoniques
$$
\varinjlim_i H^n (X_i, F_i) \to H^n (X, F_{\infty})
$$
sont des isomorphismes.
\end{enonce*}

Compte\pageoriginale tenu de \hyperref[VII.lem5.6]{5.6}, cela résulte de la conjonction de
\hyperrefext[VI][VI.thm8.7.3]{8.7.3} et \hyperrefext[VI][VI.prop8.5.2]{8.5.2}.
On appliquera surtout \hyperref[VII.thm5.7]{5.7} dans le cas particulier suivant,
explicité également dans

\begin{enonce*}{Corollaire 5.8}\etiquette[5.8]{VII.cor5.8}
Avec les hypothèses et notations précédentes pour $(X_i)_{i \in I}$,
$X$, $u_{ij}$, $u_i$, soit $i_0 \in I$, et soit $F_{i_0}$ un faisceau
abélien sur $X_{i_0}$. Pour tout $i \geq i_0$, soit $F_i = u_{i_0 i}
{}^{\ast}(F_{i_0})$, soit de plus $F_{\infty} = u^{\ast}_i
(F_{i_0})$. Sous ces conditions, les homomorphismes canoniques
$$
\varinjlim_i H^n (X_i, F_i) \to H^n (X, F_{\infty})
$$
sont des isomorphismes.
\end{enonce*}

Voici cependant un cas parfois utile qui relève de \hyperref[VII.thm5.7]{5.7} et non de \hyperref[VII.cor5.8]{5.8}:

\begin{enonce*}{Corollaire 5.9}\etiquette[5.9]{VII.cor5.9}
Avec les hypothèses et notations de \hyperref[VII.thm5.7]{5.7} pour $(X_i)_{i \in I}$, $X$,
$u_{ij}$, $u_i$, soit $G_{i_0}$ un schéma en groupes commutatifs
localement de présentation finie sur $X_{i_0}$ (où $i_0 \in I$ est
donné). Pour tout $i \geq i_0$, soit $G_i = \fprod{G_{i_0}}{X_i}{X_{i_0}}$,
et soit de même $G_{\infty} = \fprod{G_{i_0}}{X}{X_{i_0}}$. Alors
les homomorphismes canoniques
$$
\varinjlim_i H^n (X_i, G_i) \to H^n (X, G_{\infty})
$$
sont des isomorphismes.
\end{enonce*}

En effet, nous savons que $G_{i_0}$ définit un faisceau sur
$(\Sch)_{/X_{i_0}}$ (\cf 2 (a)), et supposant que $i_0$ est un plus
petit objet pour $I$ (ce qui est loisible), d'où un foncteur
canonique $\Gg \to (\Sch)_{/X_{i_0}}$, on trouve par composition un
foncteur contravariant $F$ sur $\Gg$, dont la restriction à $\Gg_i$
est\pageoriginale canoniquement isomorphe au faisceau sur $X_i$
défini par $G_i$. En particulier $F$ est un faisceau sur $\Gg$, et
nous pouvons lui appliquer \hyperref[VII.thm5.7]{5.7}. L'hypothèse que
$G_{i_0}$ est localement de présentation finie sur $X_{i_0}$ sert à
assurer que le faisceau $F_{\infty}$ de 5,7 est isomorphe, par
l'homomorphisme naturel $F_{\infty} \to \faisc(G_{\infty})$, au
faisceau $\faisc(G_{\infty})$ défini par $G_{\infty}$ (comme il
résulte en effet aussitôt de \EGA IV 8.8.2 (i), en regardant les
valeurs des faisceaux envisagés sur les objets du site étale
restreint de $G_{\infty}$ et appliquant \hyperref[VII.lem5.6]{5.6}). Cela
achève de prouver \hyperref[VII.cor5.9]{5.9}.

\begin{enonce*}{Corollaire 5.10}\etiquette[5.10]{VII.cor5.10}
Avec les hypothèses et notations de \hyperref[VII.thm5.7]{5.7} pour $(X_i)_{i \in I}$, $X$,
$u_{ij}$, $u_i$, soit $F$ un faisceau en groupes commutatifs
 sur $X$,
alors les homomorphismes canoniques
$$
\varinjlim_{i} H^n (X_i, u_{i\ast} (F)) \to H^n (X, F)
$$
sont des isomorphismes.
\end{enonce*}

Les énoncés \hyperref[VII.thm5.7]{5.7} à \hyperref[VII.cor5.10]{5.10} se généralisent en des énoncés correspondants
pour les $\Ext^i$ de faisceaux de Modules, grâce à \hyperrefext[VI][VI.cor8.7.9]{8.7.9}, que nous
nous dispensons de répéter dans le cas particulier présent. Nous
allons par contre expliciter, pour références ultérieures, les
variantes de \hyperref[VII.cor5.8]{5.8} et 5,9 en termes de foncteurs $R^n f_{\ast}$.

\begin{enonce*}{Corollaire 5.11}\etiquette[5.11]{VII.cor5.11}
Soient $I$ un ensemble préordonné filtrant croissant, $(X_i)_{i \in
I}$ et $(Y_i)_{i \in I}$ deux systèmes projectifs de schémas, à
morphismes de transition affines, $X = \varprojlim_{i} X_i$, $Y =
\varprojlim_{i} Y_i$, $(f_i)_{i \in I}$ un système projectif de
morphismes $f_i : X_i \to Y_i$ quasi-compacts et quasi-séparés, $i_0
\in I$, $F_{i_0}$ un faisceau sur $X_{i_0}$. Pour tout $i \geq i_0$,
soit $F_i$ le faisceau\pageoriginale sur $X_i$ image inverse de
$F_{i_0}$. Soit de même $f : X \to Y$ déduit de $(f_i)$ et $F$ le
faisceau sur $X$ image inverse de $F_{i_0}$. Sous ces conditions,
l'homomorphisme canonique
$$
\varinjlim_{i} u^{\ast}_i (R^n f_{i\ast} (F_i)) \to R^n
f_{\ast} (F)
$$
est un isomorphisme (où $u_i : X \to X_i$ est le morphisme canonique).
\end{enonce*}

On peut supposer que $i_0$ est un plus petit objet de $I$, et la
question étant locale sur $Y_{i_0}$, que $Y_{i_0}$ est affine, de
sorte que les $Y_i$ et les $X_i$ sont quasi-compacts et
quasi-séparés. On est alors sous les conditions d'application de
\hyperrefext[VI][VI.thm8.7.3]{8.7.3} (où on se ramène à \hyperref[VII.cor5.8]{5.8}, compte
tenu du calcul des $R^{n} f_{\ast}$ comme faisceaux associés à des
préfaisceaux, On prouve de même :

\begin{enonce*}{Corollaire 5.12}\etiquette[5.12]{VII.cor5.12}
Même énoncé que \hyperref[VII.cor5.11]{5.11}, à cela près que maintenant $F_{i_0}$ désigne un
schéma en groupes commutatifs localement de présentation finie sur
$X_{i_0}$, et $F_i$, $F$ ses images réciproques (au sens du changement
de base pour les schémas).
\end{enonce*}

On fera attention que \hyperref[VII.cor5.9]{5.9} (\resp \hyperref[VII.cor5.12]{5.12}) n'est pas un cas particulier de
\hyperref[VII.cor5.8]{5.8} (\resp \hyperref[VII.cor5.11]{5.11}), \cf \hyperref[VII.sec2]{2} a) remarque \hyperref[VII.rem2.1]{2.1}.

\begin{enonce*}{Corollaire 5.13}\etiquette[5.13]{VII.cor5.13}
Avec les notations de \hyperref[VII.cor5.11]{5.11}, soit $F$ un faisceau en groupes abéliens
sur $X$. Alors les homomorphismes canoniques
$$
\varinjlim_i v_i^{\ast} (R^n f_{i\ast} (u_{i\ast} (F))) \to
R^n f_{\ast} (F)
$$
sont des isomorphismes, (où $u_i : X \to X_i$ et $v_i : Y \to Y_i$
sont les homomorphismes canoniques).
\end{enonce*}

La démonstration est essentiellement la même que précédemment.

\begin{enonce*}[remark]{Remarques 5.14}\etiquette[5.14]{VII.rem5.14}
\begin{enumerate}
\renewcommand{\theenumi}{\alph{enumi}}
\renewcommand{\labelenumi}{\theenumi)}
\item Les\pageoriginale résultats de passage à la limite précédents
  sont valables
pour le $H^{0}$ \resp les $f_{\ast}$, pour des faisceaux
\emph{d'ensembles} (au lieu de faisceaux abéliens), et pour le $H^{1}$
\resp les $R^1 f_{\ast}$ pour des faisceaux de groupes (non
nécessairement commutatifs), - où les $H^1$ et $R^1 f_{\ast}$ non
commutatifs sont définis comme d'habitude en termes de torseurs (ou
fibrés principaux homogènes) (\cf thèse de Giraud \footnote{citée p.\pagehyperref[page7]{7}
note de bas de page $(\ast)$}). Cela peut se vérifier directement dans
le contexte général de Il est certainement possible (et sans doute utile) de
donner également une variante pour les $H^2$ non commutatifs de « liens  », étudiés par Giraud.

\item Les résultats de passage à la limite pour les sites fibrés
développés dans VI supposaient seulement que la catégorie
base était
une catégorie filtrante sans qu'il ait été nécessaire
de supposer
qu'elle soit associée à un ensemble préordonné
filtrant. Le cas d'une
catégorie d'indices filtrante arbitraire est également le cadre
naturel pour les énoncés développés dans le
présent ${\rm n^{\circ}}$. Si nous
nous sommes placés dans un cadre trop restrictif, cela était pour
pouvoir donner des références correctes à \EGA IV, où
l'on suppose
malencontreusement dans les questions de passage à la limite (à partir
du par. 8) que la catégorie d'indices est définie en termes d'un
ensemble préordonné filtrant. Nous admettrons cependant par la suite
que tous les résultats utilisés de \EGA IV sont valables pour des
catégories d'indices filtrantes quelconques, (les démonstrations
données dans \loccit étant valables, essentiellement sans changement,
dans ce cas plus général). Aussi, nous utiliserons également sans
autre commentaire les résultats du présent ${\rm n^{\circ}}$
dans le cas où $I$
est remplacé par une catégorie d'indices filtrante arbitraire.
\end{enumerate}
\end{enonce*}

\ifx\danslelivre\undefined
\end{document}
\fi
